% ============================================================================
%  97-doc-them.tex — tài liệu tham khảo có chú giải và hướng đọc tiếp
%  Ngày tạo: 2026-09-07 | Sửa gần nhất: 2026-09-07
% ============================================================================
\documentclass[english.tex]{subfiles}
\begin{document}
\chapter*{Tài liệu tham khảo có chú giải}
\ptrtarget{Z/97}
\addcontentsline{toc}{chapter}{Tài liệu tham khảo có chú giải}
\markboth{Tài liệu tham khảo có chú giải}{}

\begin{tomtat}
Danh sách \emph{Nguồn} ở cuối sách liệt kê đầy đủ mọi tài liệu được trích dẫn. Phụ lục này làm việc khác: nó nhóm chúng theo \emph{mục đích sử dụng}, nói rõ mỗi tài liệu dùng để làm gì, đọc lúc nào, và đọc bao nhiêu là đủ. Nếu chỉ nhớ một điều: chọn \emph{một} tài liệu cho nhánh đang yếu nhất của bạn và dùng nó đến nơi đến chốn --- đọc song song bốn quyển về viết là cách chắc chắn nhất để không áp dụng được quyển nào.
\end{tomtat}

\section*{A. Nếu chỉ đọc hai thứ}

\emph{Gopen và Swan, ``The Science of Scientific Writing''}\cite{gopenswan1990}. Mười trang, và là mười trang có tỉ lệ lợi ích trên công sức cao nhất trong toàn bộ danh sách này. Nó không dạy luật viết mà đổi câu hỏi: người đọc chờ gì ở vị trí nào. Toàn bộ \ptr{C/12} dựa trên bài này. Đọc chậm, và làm thử các ví dụ sửa của họ thay vì chỉ đọc lướt.

\emph{Williams, \emph{Style: Lessons in Clarity and Grace}}\cite{williams2016}. Quyển cốt lõi về viết rõ ở cấp câu. Điểm mạnh riêng: mỗi nguyên tắc đi kèm hàng chục cặp câu trước--sau, nên bạn học bằng ví dụ chứ không bằng lời khuyên. Đọc năm chương đầu; phần còn lại tra khi cần. Nền của \ptr{C/11}.

\section*{B. Từ vựng và cách học từ}

\emph{Nation, \emph{Learning Vocabulary in Another Language}}\cite{nation2013}: tham chiếu đầy đủ về việc học từ vựng, kể cả khung bốn nhánh ở \ptr{E/19}. Dày và mang tính học thuật; dùng để tra chứ không đọc tuần tự. Nếu bạn muốn bản ngắn và thực hành hơn, lấy\cite{webbnation2017}.

\emph{Schmitt, \emph{Vocabulary in Language Teaching}}\cite{schmitt2000}: nhẹ hơn, tập trung vào mặt tâm lý học của việc học từ --- vì sao một số từ dính và một số thì không.

\emph{Coxhead, danh sách từ học thuật}\cite{coxhead2000} và \emph{danh sách theo tần suất}\cite{ngsl}: không phải sách để đọc mà là danh sách để đối chiếu. Dùng một lần để biết mình đang thiếu tầng nào, rồi bỏ sang một bên.

\emph{Chung và Nation về từ kỹ thuật}\cite{chungnation2003}: bài báo ngắn, nền của \ptr{B/10}. Đọc nếu bạn muốn hiểu vì sao từ bán kỹ thuật lại nguy hiểm hơn từ hoàn toàn lạ.

\section*{C. Viết}

\emph{Swales và Feak, \emph{Academic Writing for Graduate Students}}\cite{swalesfeak2012}: sách bài tập theo thể loại, và là quyển sát nhất với việc viết bài báo. Dùng như một khoá luyện: làm bài tập chứ đừng đọc suông. Mô hình CARS ở \ptr{C/13} đến từ công trình gốc của Swales\cite{swales1990}, quyển học thuật hơn và chỉ nên đọc nếu bạn quan tâm tới phân tích thể loại.

\emph{Booth và cộng sự, \emph{The Craft of Research}}\cite{booth2016}: về cách dựng một lập luận nghiên cứu và trình bày nó. Mạnh ở chỗ nối việc \emph{đặt câu hỏi nghiên cứu} với việc viết --- thứ mà các sách văn phong bỏ qua.

\emph{Zinsser, \emph{On Writing Well}}\cite{zinsser2006} và \emph{Pinker, \emph{The Sense of Style}}\cite{pinker2014}: hai quyển về văn phong nói chung. Zinsser ngắn và thực dụng; Pinker giải thích \emph{vì sao} các nguyên tắc lại hoạt động, dựa trên tâm lý học ngôn ngữ --- bao gồm cả lời nguyền tri thức ở \ptr{C/15}.

\emph{Strunk và White, \emph{The Elements of Style}}\cite{strunkwhite1999}: rất ngắn và rất nổi tiếng. Đọc như một tài liệu lịch sử và một danh sách nhắc, không như luật --- một số quy tắc trong đó bị các nhà ngôn ngữ học phản bác có cơ sở.

\emph{Sword, \emph{Stylish Academic Writing}}\cite{sword2012}: bằng chứng rằng văn học thuật không bắt buộc phải nặng nề, kèm phân tích các bài viết rõ trong nhiều ngành. Đọc khi bạn nghi ngờ rằng ``viết rõ sẽ làm bài nghe kém học thuật''.

\emph{Graff và Birkenstein, \emph{They Say / I Say}}\cite{graff2021}: các khuôn câu để định vị ý của bạn so với người khác. Hữu ích đúng ở nước đi 2 của CARS. Dùng như ngân hàng cụm, không đọc tuần tự.

\emph{Alley, \emph{The Craft of Scientific Writing}}\cite{alley2018}: mạnh về phần trình bày kỹ thuật --- hình, bảng, cấu trúc phần kết quả.

\emph{Silvia, \emph{How to Write a Lot}}\cite{silvia2007}: không phải sách về văn phong mà về \emph{thói quen viết}. Ngắn, và giải quyết đúng vấn đề của \ptr{C/15} và \ptr{E/21}: làm sao viết đều đặn khi bận.

\section*{D. Đọc và hiểu sâu}

\emph{Keshav, ``How to Read a Paper''}\cite{keshav2007}: hai trang, phương pháp ba lượt ở \ptr{B/05} và \ptr{B/07}. Đọc trong mười phút, dùng cả đời.

\emph{Adler và Van Doren, \emph{How to Read a Book}}\cite{adler1972}: bốn tầng đọc, viết từ 1940 và vẫn dùng được. Dài; đọc phần về đọc phân tích và bỏ phần danh mục sách kinh điển.

\emph{Toulmin, \emph{The Uses of Argument}}\cite{toulmin1958}: mô hình lập luận ở \ptr{B/06}. Học thuật; đủ để biết mô hình qua chương này, chỉ đọc bản gốc nếu bạn quan tâm tới triết học lập luận.

\emph{Hyland về siêu diễn ngôn}\cite{hyland2005} \emph{và về rào đón}\cite{hyland1998}: nền của \ptr{B/08} và \ptr{C/14}. Đây là chỗ để đọc nếu bạn muốn hiểu kỹ cách tác giả báo hiệu mức cam kết.

\emph{Grice về hàm ý}\cite{grice1975}: bài báo kinh điển về những gì người nói ngụ ý mà không nói ra. Ngắn và khó; đọc chậm, và nó thay đổi cách bạn nghe người khác nói.

\section*{E. Ngôn từ: sắc thái, kết hợp, ẩn dụ}

\emph{Sinclair, \emph{Corpus, Concordance, Collocation}}\cite{sinclair1991}: nguồn của ý ``ngôn ngữ được tạo theo hai nguyên tắc'' ở \ptr{D/17}. Học thuật nhưng đọc được.

\emph{Lewis, \emph{The Lexical Approach}}\cite{lewis1993}: lập luận đặt cụm vào trung tâm việc dạy và học thay vì ngữ pháp. Đọc nếu bạn muốn thiết kế lại cách học của mình quanh cụm.

\emph{Wray, \emph{Formulaic Language}}\cite{wray2002} \emph{và Hoey, \emph{Lexical Priming}}\cite{hoey2005}: hai quyển đi sâu hơn về cùng hiện tượng. Chỉ đọc khi bạn đã áp dụng xong hai quyển trên.

\emph{Lakoff và Johnson, \emph{Metaphors We Live By}}\cite{lakoff1980}: nền của \ptr{D/18}. Dễ đọc một cách bất ngờ với một quyển có ảnh hưởng lớn đến vậy. Đọc phần đầu; phần sau về triết học có thể bỏ.

\emph{Biber và cộng sự, ngữ pháp dựa ngữ liệu}\cite{biber1999}: không đọc tuần tự. Dùng để tra khi bạn muốn biết một cấu trúc phổ biến ra sao trong văn học thuật so với văn nói.

\emph{Louw về màu sắc đánh giá}\cite{louw1993}: bài viết nguồn của hiện tượng ở \ptr{D/16} mục 2.

\section*{F. Ngữ pháp và tra cứu}

\emph{Swan, \emph{Practical English Usage}}\cite{swan2016}: dùng như từ điển ngữ pháp --- tra đúng mục khi vướng, không đọc tuần tự. Đây là quyển tra cứu đáng có nhất.

\emph{Murphy, \emph{English Grammar in Use}}\cite{murphy2019}: sách bài tập theo mục. Dùng khi bạn phát hiện mình mắc lặp lại một lỗi cụ thể (mạo từ, thì, giới từ) --- làm đúng mục đó rồi quay lại đọc.

\emph{Quirk và cộng sự}\cite{quirk1985}: ngữ pháp tham chiếu toàn diện, gần hai nghìn trang. Chỉ tra, và chỉ khi các nguồn trên không trả lời được.

\section*{G. Công cụ trực tuyến}

\emph{Từ điển cho người học}\cite{oald,ldoce}: từ điển mặc định của bạn. Đọc cả ba mục nhãn văn vực, ví dụ, và cụm đi kèm.

\emph{Ngữ liệu}\cite{coca}: để kiểm một tổ hợp cụ thể, không để đọc lan man.

\emph{Từ điển gốc từ}\cite{etymonline}: đoán nghĩa và văn vực qua gốc từ (\ptr{A/04}, \ptr{D/16}).

\emph{Ngân hàng cụm học thuật}\cite{manchesterphrasebank}: mẫu câu xếp theo chức năng. Lấy \emph{cấu trúc} của nó rồi điền bằng cụm từ chính ngành bạn.

\emph{Từ điển tổng quát}\cite{mw} \emph{và tài liệu viết học thuật trực tuyến}\cite{purdueowl}: dùng để tra định dạng trích dẫn và các quy ước trình bày.

\section*{H. Nền tâm lý học của việc luyện tập}

\emph{Roediger và Karpicke về hiệu ứng kiểm tra}\cite{roediger2006} \emph{và Cepeda và cộng sự về ôn giãn cách}\cite{cepeda2006}: hai bài báo là cơ sở cho phần \emph{Ôn nhanh} của quyển này và cho lịch ôn ở \ptr{E/21}. Đọc phần tóm tắt và phần thảo luận là đủ.

\emph{Ahrens về hệ ghi chú liên kết}\cite{ahrens2017}: cách tổ chức ghi chú để chúng được dùng lại. Nhẹ; lấy nguyên tắc ``một ghi chú một ý, viết bằng câu của mình'' và bỏ qua phần công cụ cụ thể.

\emph{Krashen}\cite{krashen1985}: đọc để hiểu một cực của cuộc tranh luận mà \ptr{E/19} tổng hợp. Đừng đọc nó như kết luận cuối cùng; đọc nó như một lập luận có ảnh hưởng và có giới hạn đã biết.

\emph{Day và Bamford về đọc rộng}\cite{daybamford1998}: nguyên tắc chọn tài liệu ở đúng mức cho nhánh 1.

\section*{I. Lịch sử từ vựng và các nguồn nền}

\emph{Firth}\cite{firth1957} cho câu phát biểu nền của \ptr{D/17};\cite{ogden1923} và\cite{west1953} cho nguồn gốc ý tưởng vốn từ hạn chế;\cite{cobuild1987} cho từ điển đầu tiên xây từ ngữ liệu;\cite{zipf1949} cho phân bố tần suất từ ở \ptr{A/02};\cite{halliday1976} cho liên kết văn bản ở \ptr{C/12};\cite{brownlevinson1987} cho lý thuyết lịch sự ở \ptr{C/14};\cite{hunation2000} và\cite{laufer2010} cho các nghiên cứu về ngưỡng độ phủ từ vựng;\cite{nation2006} cho ước lượng vốn từ cần thiết;\cite{fish2011} và\cite{lewis1949} cho hai góc nhìn khác về câu và về ngôn ngữ. Nhóm này là nguồn nền --- tra khi bạn muốn truy tới gốc một ý cụ thể, không phải để đọc tuần tự.
\end{document}
